\section{Methodology}
\label{sec:methodology}

\subsection{Dataset}
\label{sec:dataset}

We use the \dataset{}~\citep{jikeli2024goldstandard}, a collection of 11,311 English tweets annotated for antisemitic content according to the International Holocaust Remembrance Alliance (IHRA) Working Definition of Antisemitism~\citep{ihra2016working}. The dataset was curated by searching Twitter using four keyword groups that commonly appear in antisemitic discourse. Table~\ref{tab:keyword_dist} summarizes the keyword distribution and per-keyword antisemitism base rates.

\begin{table}[h]
\centering
\caption{Keyword distribution in the \dataset{} dataset. The antisemitism base rate varies dramatically across keyword groups, reflecting the different discourse contexts in which each term appears.}
\label{tab:keyword_dist}
\begin{tabular}{lrrr}
\toprule
\textbf{Keyword} & \textbf{Tweets} & \textbf{Antisemitic} & \textbf{Base Rate} \\
\midrule
Jews    & 6,365 & 725  & 11.4\% \\
Israel  & 4,134 & 664  & 16.1\% \\
ZioNazi &   529 & 467  & 88.3\% \\
Kikes   &   283 &  97  & 34.3\% \\
\midrule
\textbf{Total} & \textbf{11,311} & \textbf{1,953} & \textbf{17.3\%} \\
\bottomrule
\end{tabular}
\end{table}

Several observations about the keyword distribution are noteworthy. The ``Jews'' keyword group is the largest (56.3\% of tweets) but has the lowest antisemitism base rate (11.4\%), reflecting the fact that most tweets mentioning Jewish people are not antisemitic --- they discuss culture, news, religion, or history. The ``Israel'' group (36.6\%) has a somewhat higher base rate (16.1\%), capturing the contested terrain where political criticism may or may not cross into antisemitism. The ``ZioNazi'' group, despite being small (4.7\%), has an extremely high base rate (88.3\%) because the term itself --- equating Zionism with Nazism --- is considered antisemitic under the IHRA definition through its trivialization of the Holocaust. Finally, the ``Kikes'' group (2.5\%) shows a moderate base rate (34.3\%) because many tweets containing this slur are counterspeech (people denouncing the term) rather than antisemitic usage.

The dataset spans from January 2019 to April 2023, covering significant geopolitical events including the 2021 Israel--Gaza conflict (reflected in 2,591 tweets from 2021) and the COVID-19 pandemic period (3,712 tweets from 2020, which saw a documented rise in antisemitic conspiracy theories). Tweets have a mean length of 204 characters (30.5 words), with a maximum of 968 characters.

The overall class distribution is imbalanced: 9,358 tweets (82.7\%) are labeled not antisemitic and 1,953 (17.3\%) are antisemitic, yielding an imbalance ratio of 4.79:1. This degree of imbalance is sufficient to cause classifiers trained with standard cross-entropy to develop a bias toward the majority class.

\subsubsection{Data Splitting}

We create train/dev/test splits with an 80/10/10 ratio using \emph{joint stratification} on both the binary label and the keyword column. This ensures that each split maintains the same label distribution (17.3\% positive) \emph{and} the same keyword composition, preventing the model from encountering unfamiliar keyword distributions at test time. After preprocessing (deduplication and length filtering), the resulting splits contain 8{,}479 training, 1{,}060 development, and 1{,}060 test samples (187 antisemitic). All multi-seed numbers in this paper are reported on this 1{,}060-tweet test set.

\subsubsection{Text Preprocessing}

We apply the following preprocessing pipeline to all tweets before tokenization:
\begin{enumerate}[leftmargin=*,itemsep=1pt]
    \item \textbf{Username masking}: All @-mentions are replaced with the token \texttt{@USER} to protect privacy and reduce vocabulary size without losing the signal that a user was mentioned.
    \item \textbf{URL removal}: All URLs are replaced with the token \texttt{[URL]}, as link destinations are not available and URLs are not informative for classification.
    \item \textbf{Hashtag normalization}: The \texttt{\#} symbol is removed but the hashtag text is preserved (e.g., \texttt{\#FreePalestine} $\rightarrow$ \texttt{FreePalestine}), retaining the semantic content of hashtags.
    \item \textbf{Whitespace normalization}: Multiple spaces and newlines are collapsed to single spaces.
    \item \textbf{Case preservation}: We do not lowercase text, as capitalization carries emphasis and intensity signals (e.g., ``ALL JEWS'' vs.\ ``all Jews'').
    \item \textbf{Emoji preservation}: Emojis are retained, as they can carry sentiment, sarcasm, or threatening intent.
\end{enumerate}

\subsection{Models}
\label{sec:models}

We evaluate seven model configurations spanning two paradigms: fine-tuned transformers (Models A--C) and prompted large language models (Models D1--D3).

\subsubsection{Fine-Tuned Transformers}

All fine-tuned models follow the same architecture: a pretrained transformer encoder followed by a two-class linear classification head. They differ in the pretrained backbone, the loss function, and the data augmentation strategy.

\paragraph{Model A: RoBERTa-base (125M parameters).}
Our baseline fine-tuned model uses \modelA{}~\citep{liu2019roberta}, a robustly optimized BERT variant pretrained on a 160~GB corpus of English text. RoBERTa uses byte-pair encoding (BPE) tokenization with a 50,265-token vocabulary. We use weighted cross-entropy loss with class weights $[1.0, 4.79]$ reflecting the inverse class frequency, which increases the gradient contribution of antisemitic samples without oversampling.

\paragraph{Model B: DeBERTa-v3-base (184M parameters).}
\modelB{}~\citep{he2023debertav3} introduces two architectural innovations over BERT: (i)~\emph{disentangled attention}, which uses separate vectors for content and position rather than combining them into a single embedding, enabling more expressive attention patterns; and (ii)~ELECTRA-style pretraining with replaced token detection, which provides more efficient training signal than masked language modeling. We hypothesize that disentangled attention will be particularly helpful for antisemitism detection, where the position of identity terms relative to predicates and modifiers determines meaning. We use the same weighted cross-entropy loss as Model~A.

\paragraph{Model B2: HateBERT (110M parameters).}
\modelBB{}~\citep{caselli2021hatebert} is a BERT-base model further pretrained on approximately 1.4~million posts from banned or quarantined hateful Reddit communities. This domain-specific pretraining exposes the model to the vocabulary, rhetorical structures, and coded language of online hate speech. However, since HateBERT's pretraining corpus is Reddit-based while our target domain is Twitter, there may be a domain mismatch in terms of text length, style, and community norms. We include HateBERT to test whether domain-specific pretraining on hateful content provides an advantage over architecturally superior but general-purpose models like DeBERTa-v3.

\paragraph{Model C: DeBERTa-v3 + Focal Loss + Keyword Masking (184M parameters).}
Our proposed model combines three strategies:
\begin{itemize}[leftmargin=*,itemsep=2pt]
    \item \textbf{DeBERTa-v3-base} as the backbone, for its superior contextual understanding.
    \item \textbf{Focal loss}~\citep{lin2017focal} instead of weighted cross-entropy. Focal loss applies a modulating factor $(1 - p_t)^\gamma$ to the cross-entropy loss, where $p_t$ is the predicted probability of the correct class. With $\gamma = 2.0$, the loss is down-weighted for well-classified examples and focused on hard, ambiguous boundary cases. We set $\alpha = 0.75$ to further increase the weight of the positive (antisemitic) class. This differs from class weighting in that it focuses training not just on the minority class but specifically on examples that the model finds difficult --- which in antisemitism detection are often the nuanced cases near the criticism--hate boundary.
    \item \textbf{Keyword masking augmentation} as counterfactual regularization. During training, for each input with probability $p_{\text{mask}} = 0.3$, all identity-group keywords (``jews,'' ``jewish,'' ``israel,'' ``israeli,'' ``zionazi,'' ``zionist,'' ``kike,'' ``kikes,'' and morphological variants) are replaced with the \texttt{[MASK]} token. The label is preserved. This forces the model to classify the tweet based on \emph{context} --- rhetorical structure, predicates, conspiratorial framing, dehumanizing language --- rather than the mere presence of identity terms. At inference time, no masking is applied.
\end{itemize}

The keyword masking probability $p_{\text{mask}} = 0.3$ was chosen to balance two objectives: enough masking to break the keyword--label correlation (we want the model to see both masked and unmasked versions), but not so much that the model never sees the keywords and cannot learn their interaction with context. In implementation, the mask token is resolved from the tokenizer (\texttt{tokenizer.mask\_token}) so the augmentation works equally for DeBERTa's \texttt{[MASK]} and RoBERTa's \texttt{<mask>}.

\paragraph{Multi-architecture variants.}
To test whether the counterfactual-invariance effect we study below transfers across scale and family, we additionally train Model~C and CCI~v2 on two larger backbones from different architectural families: DeBERTa-v3-large (440M parameters, same family) and RoBERTa-large (355M parameters, standard attention) --- six seeds per cell. Hyperparameters are halved-batch / doubled-grad-accumulation copies of the base recipe (effective batch size unchanged at 16), with learning rate $1 \times 10^{-5}$ per the original DeBERTa-v3 and RoBERTa fine-tuning recipes.

\paragraph{Counterfactual Logit Pairing baselines (Davani CLP).}
For comparison against the closest published prior art we train Counterfactual Logit Pairing baselines~\citep{garg2019counterfactual,davani2021improving} at four $\lambda$ settings ($\lambda \in \{0.5, 1.0, 2.0\}$ with the heuristic symmetry filter, plus $\lambda = 1.0$ without the filter to isolate the filter's contribution). All Davani baselines share the Model~C backbone and recipe; they differ only in the auxiliary pair-loss term, applied in logit space with L\textsubscript{1} divergence.

\subsubsection{Prompted Large Language Models}

We evaluate four prompted LLM families --- Mistral-7B-Instruct~\citep{jiang2023mistral}, Qwen-2.5-7B-Instruct~\citep{qwen25}, Llama-3.1-8B-Instruct~\citep{llama31}, and Llama-3.3-70B (the latter via the Groq inference API) --- using three prompting strategies. The 7B / 8B models use 4-bit NormalFloat (NF4) quantization via bitsandbytes~\citep{dettmers2022gptint} and run on a single A100; the 70B model is queried over the Groq API in fp16. Temperature is set to 0 (greedy decoding) for reproducibility. Chat templates are resolved automatically from each tokenizer.

\paragraph{Model D1: Zero-Shot.}
The prompt provides the IHRA definition, explains important distinctions (criticism vs.\ antisemitism, counterspeech vs.\ hate), and asks the model to classify a single tweet as ``antisemitic'' or ``not antisemitic.'' No examples are provided.

\paragraph{Model D2: Few-Shot (5 examples).}
The prompt includes five carefully selected demonstration examples: two antisemitic (conspiracy theory, Holocaust trivialization), two not antisemitic (political criticism, cultural discussion), and one edge case (counterspeech quoting a slur to denounce it). Examples were selected for diversity across keyword groups and difficulty levels.

\paragraph{Model D3: Guided Chain-of-Thought.}
Following~\citet{patel2025evaluating}, we structure the model's reasoning into four explicit steps: (1)~identify the target, (2)~check for antisemitic tropes (conspiracy theories, dual loyalty, Holocaust trivialization, dehumanization), (3)~check for legitimate speech (factual reporting, policy criticism, counterspeech), and (4)~classify. The model is given 300 max tokens to produce its step-by-step reasoning before the final classification.

For all LLM configurations, we parse the model's output to extract the binary classification. Responses that do not contain a clear ``antisemitic'' or ``not antisemitic'' label are marked as ``invalid'' and excluded from metric computation, with the invalid rate reported as a separate metric.

\subsection{Counterfactual Calibration Invariance (CCI v2)}
\label{sec:cci}

Beyond keyword masking (Model~C), we propose a stronger training-time intervention that targets counterfactual stability directly. Counterfactual Calibration Invariance v2 (CCI v2) generalises Counterfactual Logit Pairing~\citep{garg2019counterfactual,davani2021improving} along two axes: divergence in \emph{probability} space rather than L\textsubscript{p} in logit space, and a per-group learnable temperature.

\paragraph{Counterfactual swap engine.}
Given a tweet $x$ containing a Jewish religious-ethnic token $t \in T_{\text{Jewish}} = \{\text{jew}, \text{jewish}, \text{jews}, \ldots\}$, we generate up to $K = 5$ counterfactual variants $\{x'_1, \ldots, x'_K\}$ by substituting $t$ with a token drawn from a religiously-symmetric vocabulary $T_{\text{alt}} = \{\text{Muslim}, \text{Christian}, \text{Hindu}, \text{Buddhist}, \text{Atheist}\}$. Casing and morphology are preserved. The engine is wrapped in a \emph{heuristic symmetry classifier} that rejects swaps for which the label may not be preserved: (i) if $x$ contains an antisemitic slur (\texttt{kike}, \texttt{zionazi}, ...) the swap is rejected because the slur is identity-specific; (ii) if $x$ contains only a geopolitical proxy (\texttt{israel}, \texttt{zionist}) without a religious-ethnic anchor, the swap is rejected because no symmetric counterpart exists. On the 1{,}060-tweet test set, this yields a manifest of 2{,}117 symmetric pairs (\texttt{Jews} group: 1{,}725 pairs; \texttt{Israel} group: 392 pairs; \texttt{Kikes} and \texttt{ZioNazi} groups: 0 by construction).

\paragraph{Per-group temperature head.}
We attach a small head $\tau_\theta : \mathcal{G} \to \mathbb{R}_{>0}$ that maps each example's identity-group id $g \in \mathcal{G}$ = \{Jews, Israel, Kikes, ZioNazi, none\} to a scalar temperature $T_g = \exp(\log T_g)$, clipped to $[T_{\min}, T_{\max}] = [0.5, 5.0]$. The head is jointly trained with the model (added as a parameter group to the same AdamW optimizer, without weight decay on $\log T_g$). At inference we discard $\tau_\theta$ --- it modulates the loss, not the prediction.

\paragraph{Loss.}
Let $z(x) = \mathrm{logit}_1(x) - \mathrm{logit}_0(x)$ be the positive-class logit difference, $\sigma(\cdot)$ the sigmoid, and $\mathrm{JS}_b(p, q)$ the binary Jensen-Shannon divergence between Bernoulli$(p)$ and Bernoulli$(q)$. For a training mini-batch $\{(x_i, y_i, \{x'_{i,k}\}_{k=1}^{K_i})\}_{i=1}^{B}$ with per-example group ids $g_i$,
\begin{equation}
\mathcal{L}_{\text{CCI}}(x_i)
\;=\;
\mathrm{FocalLoss}(f(x_i), y_i)
\;+\;
\lambda \,\frac{1}{K_i}\sum_{k=1}^{K_i}
\mathrm{JS}_b\!\Big(\sigma\!\big(z(x_i) / T_{g_i}\big),\,
                     \sigma\!\big(z(x'_{i,k}) / T_{g_i}\big)\Big).
\label{eq:cci-loss}
\end{equation}
We set $\lambda = 1.0$. The four CCI v2 ablation cells vary two binary axes: $T_g$ \emph{learnable} vs.\ \emph{fixed} at 1.0, and divergence \emph{JS} (eq.~\ref{eq:cci-loss}) vs.\ \emph{L\textsubscript{1}} ($|\sigma(z(x_i)/T_{g_i}) - \sigma(z(x'_{i,k})/T_{g_i})|$). The fixed-T cell isolates the divergence contribution; the learned-T cell tests whether per-group calibration provides additional headroom.

\paragraph{Structural lemma: scalar Temperature Scaling cannot reduce CFR\textsubscript{hard}.}
\label{para:ts-cfr}
A standard post-hoc calibration fix --- Temperature Scaling~\citep{guo2017calibration} --- divides all logits by a single scalar $T \in \mathbb{R}_{>0}$ fitted on a held-out dev set. The following observation rules out scalar TS as a tool for reducing the hard counterfactual flip rate.

\begin{lemma}[Scalar TS preserves CFR\textsubscript{hard} at threshold $1/2$]
\label{thm:ts-cfr}
Let $f : \mathcal{X} \to \mathbb{R}^2$ be any binary classifier with positive-class logit difference $z(x) = f_1(x) - f_0(x)$, and let $\tilde z(x) = z(x) / T$ for any $T > 0$. For any counterfactual pair $(x, x')$,
\begin{equation*}
\mathbb{1}[\sigma(z(x)) \ge \tfrac{1}{2}] \;\ne\; \mathbb{1}[\sigma(z(x')) \ge \tfrac{1}{2}]
\;\;\Longleftrightarrow\;\;
\mathbb{1}[\sigma(\tilde z(x)) \ge \tfrac{1}{2}] \;\ne\; \mathbb{1}[\sigma(\tilde z(x')) \ge \tfrac{1}{2}].
\end{equation*}
\end{lemma}
\begin{proof}[Proof.] $\sigma(\cdot) \ge \tfrac{1}{2} \iff (\cdot) \ge 0$; dividing by $T > 0$ preserves sign. \qed
\end{proof}

Lemma~\ref{thm:ts-cfr} is a one-line observation about monotone rescaling. Its content is the immediate corollary: reducing CFR\textsubscript{hard} requires a method that changes the relative position of $z(x)$ and $z(x')$ --- this can be a training-time intervention (like CCI v2) \emph{or} a non-monotone post-hoc method (like multicalibration, which we show in~\S\ref{sec:results} reduces CFR\textsubscript{hard} by $8\%$ but raises CFR\textsubscript{soft, L\textsubscript{1}} by $96\%$). The lemma is not by itself a contribution; we report it as a regression-test sanity check, verified bit-exactly on all six Model~C seeds (Appendix~E).

\subsection{Training Procedure}
\label{sec:training}

All fine-tuned models are trained with the following protocol:
\begin{itemize}[leftmargin=*,itemsep=1pt]
    \item \textbf{Optimizer}: AdamW~\citep{loshchilov2019adamw} with learning rate $2 \times 10^{-5}$ and weight decay 0.01. Bias and LayerNorm parameters are excluded from weight decay.
    \item \textbf{Schedule}: Linear warmup over the first 10\% of training steps, followed by linear decay to zero.
    \item \textbf{Batch size}: 8 per GPU with gradient accumulation over 2 steps, yielding an effective batch size of 16.
    \item \textbf{Precision}: Mixed precision (fp16) training using PyTorch AMP.
    \item \textbf{Gradient clipping}: Maximum gradient norm of 1.0.
    \item \textbf{Early stopping}: Training halts when the development set \fone{} does not improve for 3 consecutive epochs. The checkpoint with the best development \fone{} is used for final evaluation.
    \item \textbf{Epochs}: Up to 5 epochs for Models A, B, B2; up to 7 epochs for Model C (focal loss converges more slowly).
    \item \textbf{Sequence length}: Maximum 128 tokens (sufficient for 99\%+ of preprocessed tweets).
    \item \textbf{Reproducibility}: All random seeds are set to 42 (Python, NumPy, PyTorch, CUDA).
\end{itemize}

\subsection{Evaluation Protocol}
\label{sec:eval_protocol}

We evaluate all models on the held-out test set (1,060 samples) using the following metrics and analyses:

\subsubsection{Classification Metrics}

\begin{itemize}[leftmargin=*,itemsep=1pt]
    \item \textbf{Macro-F1} (primary metric): The unweighted average of per-class F1 scores, treating the majority and minority classes equally. This is the standard metric for imbalanced binary classification.
    \item \textbf{Precision and Recall}: Computed for the positive (antisemitic) class, capturing the trade-off between false alarms and missed detections.
    \item \textbf{PR-AUC}: Area under the precision--recall curve, computed from predicted probabilities. This measures ranking quality and is less sensitive to threshold choice than F1.
    \item \textbf{Expected Calibration Error (ECE)}~\citep{guo2017calibration}: Measures how well predicted probabilities match observed frequencies, computed over 15 equal-width bins (the larger bin count of \citet{kumar2019verified}, used in our group-conditional protocol for stability on the smaller keyword sub-populations). A well-calibrated model with ECE $\approx 0$ is essential if predicted probabilities will be used for downstream decision-making (e.g., prioritizing content for human review).
\end{itemize}

\subsubsection{Counterfactual robustness}

We measure counterfactual stability on the 2{,}117-pair test manifest produced by the swap engine of \S\ref{sec:cci}. For a model with positive-class probability $p(x) = \sigma(z(x))$ and a pair $(x, x')$:
\begin{itemize}[leftmargin=*,itemsep=1pt]
    \item \textbf{CFR\textsubscript{hard}}: the fraction of pairs where the \emph{binary} prediction flips at threshold $0.5$. This is the metric Lemma~\ref{thm:ts-cfr} characterises.
    \item \textbf{CFR\textsubscript{soft, L\textsubscript{1}}}: $|p(x) - p(x')|$, averaged over pairs. More sensitive than CFR\textsubscript{hard} because it captures probability shifts that do not cross the threshold.
    \item \textbf{CFR\textsubscript{soft, JS}}: binary Jensen-Shannon divergence between $\sigma(z(x))$ and $\sigma(z(x'))$, averaged over pairs. The probabilistic counterpart of CFR\textsubscript{soft, L\textsubscript{1}}; we report both because the L\textsubscript{1} and JS metrics can diverge under post-hoc multicalibration (\S\ref{sec:results}).
\end{itemize}

\subsubsection{Group-conditional calibration and fairness}

\begin{itemize}[leftmargin=*,itemsep=1pt]
    \item \textbf{Per-keyword macro-F\textsubscript{1}}: macro-F\textsubscript{1} reported separately for each of the four GoldStandard2024 keyword groups.
    \item \textbf{Group-conditional ECE}: 15-bin ECE computed within each keyword group; we report the global ECE, the worst-group ECE, and the gap (worst $-$ global). A model can be well-calibrated globally and severely miscalibrated on the smaller groups; the worst-group ECE surfaces this.
    \item \textbf{FPED / FNED}~\citep{dixon2018measuring}: the spread of group-conditional false-positive and false-negative rates across the four keyword groups, computed as the sum of absolute deviations from the dataset-wide rate.
    \item \textbf{Cross-dataset transfer}: zero-shot evaluation of our best fine-tuned model on the Jewish-target subset of HateXplain~\citep{mathew2021hatexplain} and the Jewish-group subset of ToxiGen~\citep{hartvigsen2022toxigen}.
\end{itemize}

\subsubsection{Statistical significance protocol}
\label{sec:sig_protocol}

Because we evaluate along three connected axes against multiple baselines, multiple-comparison correction is essential. We compute paired-bootstrap p-values~\citep{dror2018hitchhikers} over six random seeds with $n_{\text{resample}} = 10{,}000$, alongside the exact two-sided sign-test p-value (which provides a principled lower bound on direction-only tests; with $n=6$, $p_{\text{sign}} \ge 2 \cdot (1/2)^6 = 0.03125$). Pre-specifying hypothesis families avoids the omnibus-suite pitfall flagged in~\citet{dror2018hitchhikers}.

The pivot for both families is the CCI v2 (fixed $T$, JS) cell on DeBERTa-v3-base. (We changed pivot from the originally pre-registered learned-$T$ variant after the within-cell ablation showed learned-$T$ is statistically indistinguishable from fixed-$T$ on CFR\textsubscript{hard} with $p = 0.25$, while fixed-$T$ has higher F\textsubscript{1}\textsuperscript{+} and lower FNED; we report this transparently rather than retain the strictly worse pre-registered pivot.)

\paragraph{Family 1 — Main suite (9 comparators $\times$ 6 metrics = 54 tests).}
Pivot vs.\ (a)~four base baselines (Davani CLP $\lambda{=}0.5$, Model~C focal+mask, Model~C+post-hoc-TS, Model~C+post-hoc-multicalibration), (b)~the CCI~v2 (learned $T$, JS) cell on the same base architecture (an ablation testing whether learned-$T$ adds anything), and (c)~four multi-architecture cells (Model~C and CCI v2 fixed\_js on DeBERTa-v3-large and RoBERTa-large). Metrics: F\textsubscript{1}\textsuperscript{+} is omitted from the suite because the pivot's F\textsubscript{1}\textsuperscript{+} advantage is small and we do not claim it; the six metrics tested are CFR\textsubscript{hard}, CFR\textsubscript{soft, L\textsubscript{1}}, CFR\textsubscript{soft, JS}, FPED, FNED, ECE\textsubscript{worst}.

\paragraph{Family 2 — Within-architecture transfer (2 pairs $\times$ 3 CFR metrics = 6 tests).}
Same-architecture comparisons: DeBERTa-v3-large CCI~v2 fixed\_js vs.\ DeBERTa-v3-large Model~C, and RoBERTa-large CCI~v2 fixed\_js vs.\ RoBERTa-large Model~C. These directly test whether the CFR reduction transfers \emph{at iso-architecture}, which the main suite (with its base pivot) does not test. CCI~v2 fixed\_js on DeBERTa-v3-large has $n=5$ common seeds (seed 47 was dropped during training); the paired bootstrap restricts to the intersection.

\paragraph{Correction.}
We apply Holm-Bonferroni step-down within each family at $\alpha = 0.05$ (uniformly more powerful than naive Bonferroni at the same family-wise error rate). We additionally report the naive Bonferroni-corrected threshold ($\alpha / K$ where $K$ is the family size, i.e.\ $\alpha = 0.05/54 \approx 9.3 \times 10^{-4}$ for Family 1 and $\alpha = 0.05/6 \approx 8.3 \times 10^{-3}$ for Family 2) and the exact sign-test p-value as robustness checks for the small-$n$ regime. A test is reported as significant only if both the Holm criterion and the sign-test direction-agreement criterion are met.
